\section*{अध्याय \ref{ch:b1:logic} --- तर्क, समुच्चय और प्रतिचित्रण}

\begin{solution}{exo:b1:logic:1}
निषेध, प्रत्येक परिमाणक के भीतर से $\lnot$ को धकेलते हुए
(\cref{prop:b1:logic:negquant}) और $\lnot(P \implies Q) \iff P \land \lnot
Q$ का प्रयोग करते हुए:
\begin{enumerate}
  \item $\exists x \in \R,\ \forall y \in \R,\ x + y \leq 0$;
  \item $\forall x \in \R,\ \exists y \in \R,\ xy \neq 0$;
  \item $\exists \varepsilon > 0,\ \forall \delta > 0,\ \exists x \in
        \R,\ \abs{x} \leq \delta \text{ और } \abs{f(x)} >
        \varepsilon$.
\end{enumerate}
\omterm{def:b1:logic:statement}{कथन} (1) सत्य है: $x$ दिया हो, तो $y = -x + 1$ लीजिए; तब $x + y = 1 > 0$। \omterm{def:b1:logic:statement}{कथन}
(2) सत्य है: $x = 0$ सभी $y$ के लिए $xy = 0$ को संतुष्ट करता है।
\end{solution}

\begin{solution}{exo:b1:logic:2}
सत्य सारणी, चारों स्थितियों $(P, Q)$ के लिए स/अ लिखते हुए:
\begin{center}
\begin{tabular}{cc|c|c|c|c}
$P$ & $Q$ & $P \implies Q$ & $\lnot(P \implies Q)$ & $\lnot Q$ & $P \land
\lnot Q$ \\ \hline स & स & स & अ & अ & अ \\ स & अ & अ & स & स & स \\ अ & स &
स & अ & अ & अ \\ अ & अ & स & अ & स & अ \\
\end{tabular}
\end{center}
स्तंभ $4$ और $6$ मेल खाते हैं, जिससे तुल्यता सिद्ध हो जाती है। अतः ``यदि कोई
फलन अवकलनीय है तो वह संतत है'' का निषेध है: ``ऐसा फलन है जो अवकलनीय है और
संतत नहीं'' (जो वस्तुतः असत्य \omterm{def:b1:logic:statement}{कथन} है: मूल निहितार्थ सत्य है,
\cref{ch:b1:derivative} देखिए)।
\end{solution}

\begin{solution}{exo:b1:logic:3}
\emph{प्रतिधनात्मकता।} मान लीजिए $x < 1$। तब $x^3 < 1$ (घन फलन वर्धमान है)
और $x < 1$, अतः $x^3 + x < 2$। इससे प्रतिधनात्मक \omterm{def:b1:logic:statement}{कथन} सिद्ध हो जाता है, और
उसके साथ मूल \omterm{def:b1:logic:statement}{कथन} भी।

\emph{विरोधाभास।} मान लीजिए $a > 0$ लघुतम पूर्णतः धनात्मक वास्तविक संख्या
है। तब $a/2$ पूर्णतः धनात्मक है और $a/2 < a$ ($a > 0$ के कारण), जो लघुतमता
का विरोध करता है। अतः ऐसा कोई $a$ नहीं है।
\end{solution}

\begin{solution}{exo:b1:logic:4}
\begin{enumerate}
  \item आधार स्थिति $n = 0$: $2^0 = 1 = 2^1 - 1$। पद: $n$ के लिए सर्वसमिका
  मानने पर,
        \[
        \sum_{k=0}^{n+1} 2^k = (2^{n+1} - 1) + 2^{n+1}
        = 2 \cdot 2^{n+1} - 1 = 2^{n+2} - 1 .
        \]
  \item आधार स्थिति $n = 0$: $4^0 + 5 = 6 = 3 \times 2$। पद: यदि $4^n + 5 =
  3m$, तो
        \[
        4^{n+1} + 5 = 4(4^n + 5) - 15 = 3(4m - 5),
        \]
        $3$ से विभाज्य
\end{enumerate}
\end{solution}

\begin{solution}{exo:b1:logic:5}
आगमन पद चुपचाप यह मान लेता है कि दोनों समूह (``पहली $n$'' और ``अंतिम $n$'')
अतिव्यापी हैं, जिससे साझा पेंसिलें रंग को एक समूह से दूसरे तक ले जाती हैं।
$n + 1 = 2$ के लिए दोनों समूह $\{$पहली पेंसिल$\}$ और $\{$दूसरी पेंसिल$\}$
हैं: वे असंयुक्त हैं और तर्क टूट जाता है। अतः $P(1) \implies P(2)$ कभी सिद्ध
ही नहीं हुआ और आगमन ढह जाता है --- यद्यपि $P(n) \implies P(n+1)$ प्रत्येक $n
\geq 2$ के लिए वैध है।
\end{solution}

\begin{solution}{exo:b1:logic:6}
\begin{enumerate}
  \item $x \in A \setminus B \iff x \in A \land x \notin B \iff x \in
        A \land x \in \overline{B} \iff x \in A \cap \overline{B}$.
  \item (1) और वितरणशीलता (\cref{prop:b1:logic:setalgebra}) का प्रयोग करते
  हुए: $(A \cup B) \cap \overline{C} = (A \cap \overline{C}) \cup (B \cap
  \overline{C})$।
  \item मान लीजिए $A \subseteq B$। तब $A \cup B \subseteq B$ (दोनों टुकड़े
  $B$ में हैं) और $B \subseteq A \cup B$ सदा, अतः $A \cup B = B$। मान लीजिए
  $A \cup B = B$: तब $A \cap B \subseteq A$ सदा, और $A \subseteq A \cup B =
  B$ से $A \subseteq A \cap B$ मिलता है, अतः $A \cap B = A$। मान लीजिए $A
  \cap B = A$: तब $A = A \cap B \subseteq B$। इस प्रकार तीनों शर्तें तुल्य
  हैं (हमने निहितार्थों का एक चक्र सिद्ध किया)।
\end{enumerate}
\end{solution}

\begin{solution}{exo:b1:logic:7}
\begin{enumerate}
  \item $f$ \omterm{def:b1:logic:inj}{एकैकी} है ($n + 1 = m + 1 \implies n = m$) पर \omterm{def:b1:logic:inj}{आच्छादक} नहीं: $0$
  का $\N$ में कोई \omterm{def:b1:logic:map}{पूर्वप्रतिबिंब} नहीं है।
  \item $g$ \omterm{def:b1:logic:inj}{एकैकी आच्छादक} है: $n \mapsto n - 1$ $\Z$ पर द्विपक्षीय प्रतिलोम
  है।
  \item $h$ \omterm{def:b1:logic:inj}{एकैकी} है: $\frac{x+1}{x-1} = \frac{x'+1}{x'-1}$ से $(x+1)(x'-1)
  = (x'+1)(x-1)$ मिलता है, अर्थात् $xx' - x + x' - 1 = xx' - x' + x - 1$,
  अतः $2x' = 2x$। वह $\R$ पर \omterm{def:b1:logic:inj}{आच्छादक} नहीं है: $y = \frac{x+1}{x-1}$ हल करने
  पर $x(y - 1) = y + 1$ मिलता है, जिसका $y = 1$ होने पर कोई हल नहीं है
  (समीकरण $0 = 2$ हो जाता है)। सहप्रांत $\R \setminus \{1\}$ लेने पर वही
  संगणना अद्वितीय \omterm{def:b1:logic:map}{पूर्वप्रतिबिंब} $x = \frac{y+1}{y-1}$ देती है, अतः $h
  \colon \R \setminus \{1\} \to \R \setminus \{1\}$ \omterm{def:b1:logic:inj}{एकैकी आच्छादक} है और
  $h^{-1}(y) = \frac{y+1}{y-1} = h(y)$: $h$ अपना ही प्रतिलोम है।
\end{enumerate}
\end{solution}

\begin{solution}{exo:b1:logic:8}
\begin{enumerate}
  \item $x \in f^{-1}(B \cap B') \iff f(x) \in B \cap B' \iff f(x) \in B
  \land f(x) \in B' \iff x \in f^{-1}(B) \cap f^{-1}(B')$। प्रतिबिंबों के
  लिए: $y \in f(A \cup A')$ तभी जब $A$ या $A'$ के किसी $x$ के लिए $y =
  f(x)$, अर्थात् तभी जब $y \in f(A)$ या $y \in f(A')$।
  \item यदि $y \in f(A \cap A')$, तो $y = f(x)$, जहाँ $x \in A$ और $x \in
  A'$, अतः $y \in f(A)$ और $y \in f(A')$। कठोरता: $f \colon \R \to \R$, $x
  \mapsto x^2$, $A = \{-1\}$, $A' = \{1\}$ लीजिए: तब $f(A \cap A') =
  f(\emptyset) = \emptyset$, पर $f(A) \cap f(A') = \{1\}$।
  \item ($\Leftarrow$) $A = \{x\}$ लेने पर, $x \neq x'$ के लिए $A' =
  \{x'\}$: यदि $f(x) = f(x')$, तो $f(A) \cap f(A') = \{f(x)\}$ जबकि $f(A
  \cap A') = \emptyset$, जो मानी हुई समता का विरोध करता है; अतः $f$ \omterm{def:b1:logic:inj}{एकैकी}
  है। ($\Rightarrow$) मान लीजिए $f$ \omterm{def:b1:logic:inj}{एकैकी} है और $y \in f(A) \cap f(A')$: $y
  = f(x) = f(x')$, जहाँ $x \in A$, $x' \in A'$; एकैकीयता से $x = x' \in A
  \cap A'$ मिलता है, अतः $y \in f(A \cap A')$। (2) के साथ मिलाकर समता सिद्ध
  होती है।
\end{enumerate}
\end{solution}

\begin{solution}{exo:b1:logic:9}
$g \circ f = \mathrm{id}_E$ \omterm{def:b1:logic:inj}{एकैकी} और \omterm{def:b1:logic:inj}{आच्छादक} है, अतः
\cref{prop:b1:logic:comp}~(2) से $f$ \omterm{def:b1:logic:inj}{एकैकी} है और $g$ \omterm{def:b1:logic:inj}{आच्छादक} है। उदाहरण: $E
= \N$, $F = \Z$, $f$ अंतर्विष्टि $n \mapsto n$, तथा $g \colon \Z \to \N$, $n
\geq 0$ के लिए $g(n) = n$ और $n < 0$ के लिए $g(n) = 0$। तब सभी $n \in \N$ के
लिए $g(f(n)) = n$, पर $f$ \omterm{def:b1:logic:inj}{आच्छादक} नहीं है और $g$ \omterm{def:b1:logic:inj}{एकैकी} नहीं है।
\end{solution}

\begin{solution}{exo:b1:logic:10}
$x^2 - y^2 = x - y \iff (x - y)(x + y) = x - y \iff (x - y)(x + y - 1) = 0
\iff y = x$ या $y = 1 - x$। \emph{स्वतुल्य:} $y = x$ काम करता है।
\emph{सममित:} शर्त ``$y = x$ या $y = 1 - x$'' $x$ और $y$ में सममित है (यदि
$y = 1 - x$, तो $x = 1 - y$)। \emph{संक्रमक:} मान लीजिए $x
\mathbin{\mathcal{R}} y$ और $y \mathbin{\mathcal{R}} z$; चारों स्थितियों से
गुज़रने पर $z$ हर बार $x$ या $1 - x$ के बराबर निकलता है (उदाहरणार्थ $y = 1 -
x$ और $z = 1 - y$ से $z = x$ मिलता है)। अतः $\mathcal{R}$ एक \omterm{def:b1:logic:equiv}{तुल्यता संबंध}
है और $\mathrm{cl}(x) = \{x,\, 1 - x\}$। इस वर्ग में ठीक एक अवयव तभी है जब
$x = 1 - x$, अर्थात् $x = \frac12$ के लिए।
\end{solution}

\begin{solution}{exo:b1:logic:11}
मान लीजिए $f \colon E \to \mathcal{P}(E)$ कोई \omterm{def:b1:logic:map}{प्रतिचित्रण} है और $D = \{x \in
E : x \notin f(x)\} \in \mathcal{P}(E)$ रखिए। मान लीजिए किसी $a \in E$ के
लिए $D = f(a)$। यदि $a \in D$, तो $D$ की परिभाषा से $a \notin f(a) = D$:
विरोधाभास। यदि $a \notin D$, तो $a \notin f(a)$, अतः $D$ की परिभाषा से $a
\in D$: विरोधाभास। अतः $D$ $f$ के प्रतिबिंब में नहीं है, और $f$ \omterm{def:b1:logic:inj}{आच्छादक} नहीं
है। (विशेष रूप से, कोई भी \omterm{def:b1:logic:sets}{समुच्चय} अपने \omterm{def:b1:logic:sets}{घात समुच्चय} के साथ एकैकी आच्छादन में
नहीं है: $\N$ के उपसमुच्चय पूर्णांकों से ``अधिक'' हैं।)
\end{solution}

\begin{solution}{exo:b1:logic:12}
\begin{enumerate}
  \item ($\Rightarrow$) मान लीजिए $f$ \omterm{def:b1:logic:inj}{आच्छादक} है और $\Phi(B) = \Phi(B')$। $y
  \in B$ के लिए ऐसा $x$ चुनिए कि $f(x) = y$; तब $x \in f^{-1}(B) =
  f^{-1}(B')$, अतः $y = f(x) \in B'$। अतः $B \subseteq B'$, और सममित रूप से
  $B' \subseteq B$: $\Phi$ \omterm{def:b1:logic:inj}{एकैकी} है। ($\Leftarrow$) यदि $f$ \omterm{def:b1:logic:inj}{आच्छादक} नहीं है,
  तो प्रतिबिंब के बाहर कोई $y_0 \in F$ चुनिए; तब $f^{-1}(\{y_0\}) =
  \emptyset = f^{-1}(\emptyset)$, जहाँ $\{y_0\} \neq \emptyset$, अतः $\Phi$
  \omterm{def:b1:logic:inj}{एकैकी} नहीं है।
  \item ($\Rightarrow$) मान लीजिए $f$ \omterm{def:b1:logic:inj}{एकैकी} है और $A \subseteq E$। $B =
  f(A)$ रखिए; तब $f^{-1}(B) = \{x : f(x) \in f(A)\}$, और एकैकीयता से $f(x)
  \in f(A) \iff x \in A$ मिलता है, अतः $\Phi(B) = A$: $\Phi$ \omterm{def:b1:logic:inj}{आच्छादक} है।
  ($\Leftarrow$) यदि $f$ \omterm{def:b1:logic:inj}{एकैकी} नहीं है, तो $x \neq x'$ ऐसा लीजिए कि $f(x) =
  f(x')$। प्रत्येक \omterm{def:b1:logic:map}{पूर्वप्रतिबिंब} \omterm{def:b1:logic:sets}{समुच्चय} $f^{-1}(B)$ $x$ को तभी समेटता है
  जब वह $x'$ को समेटता हो; अतः $\{x\}$ $\Phi(B)$ के रूप का नहीं है, और
  $\Phi$ \omterm{def:b1:logic:inj}{आच्छादक} नहीं है।
\end{enumerate}
\end{solution}

\begin{solution}{pb:b1:logic:1}
\textbf{1.} \emph{स्वतुल्य:} $\mathrm{id}_E$ $E$ का स्वयं पर एकैकी आच्छादन
है। \emph{सममित:} यदि $f \colon E \to F$ \omterm{def:b1:logic:inj}{एकैकी आच्छादक} है, तो
\cref{thm:b1:logic:inverse} $f^{-1} \colon F \to E$ देता है, जो स्वयं \omterm{def:b1:logic:inj}{एकैकी
आच्छादक} है। \emph{संक्रमक:} यदि $f \colon E \to F$ और $g \colon F \to G$
एकैकी आच्छादन हैं, तो \cref{prop:b1:logic:comp}~(1) कहता है कि $g \circ f
\colon E \to G$ एकैकी आच्छादन है। (यह केवल \omterm{def:b1:logic:equiv}{तुल्यता संबंध} ``की भाँति'' है:
सभी समुच्चयों का संग्रह स्वयं \omterm{def:b1:logic:sets}{समुच्चय} नहीं है --- उन्हीं विरोधाभासों के कारण
जिनकी ओर \cref{exo:b1:logic:11} संकेत करता है; महत्त्व तीनों गुणधर्मों का
है।)

\textbf{2.} यदि $f \colon E \to F$ और $g \colon F \to G$ \omterm{def:b1:logic:inj}{एकैकी} हैं, तो
\cref{prop:b1:logic:comp}~(1) से $g \circ f$ \omterm{def:b1:logic:inj}{एकैकी} है: $E \preceq G$। दूसरे
बिंदु के लिए $f$ को उसके प्रतिबिंब तक सहसीमित कीजिए: \omterm{def:b1:logic:map}{प्रतिचित्रण} $\tilde f
\colon E \to f(E)$, $x \mapsto f(x)$, $f(E)$ की रचना से \omterm{def:b1:logic:inj}{आच्छादक} है और $f$ के
\omterm{def:b1:logic:inj}{एकैकी} होने से \omterm{def:b1:logic:inj}{एकैकी}, अतः \omterm{def:b1:logic:inj}{एकैकी आच्छादक}: $E \approx f(E)$।

\textbf{3.} ($\Rightarrow$) मान लीजिए $f \colon E \to F$ \omterm{def:b1:logic:inj}{एकैकी} है और $a \in
E$ ($E \neq \emptyset$) नियत कीजिए। $s \colon F \to E$ इस प्रकार परिभाषित
कीजिए: $y \in f(E)$ होने पर $s(y)$ वह अद्वितीय $x$ है जिसके लिए $f(x) = y$
(अद्वितीयता एकैकीयता से), और अन्यथा $s(y) = a$। प्रत्येक $x \in E$ के लिए
$s(f(x)) = x$, अतः प्रत्येक $x$ प्राप्त होता है: $s$ \omterm{def:b1:logic:inj}{आच्छादक} है।
($\Leftarrow$) मान लीजिए $s \colon F \to E$ \omterm{def:b1:logic:inj}{आच्छादक} है। प्रत्येक $x \in E$
के लिए ऐसा एक $y_x \in F$ चुनिए कि $s(y_x) = x$, और $u(x) = y_x$ रखिए। यदि
$u(x) = u(x')$, तो $x = s(u(x)) = s(u(x')) = x'$: $u \colon E \to F$ \omterm{def:b1:logic:inj}{एकैकी}
है।

\textbf{4.} $n \mapsto n + 1$ $\N$ को $\N^*$ में भेजता है, \omterm{def:b1:logic:inj}{एकैकी} है ($n + 1
= m + 1 \implies n = m$) और \omterm{def:b1:logic:inj}{आच्छादक} है (प्रत्येक $m \geq 1$ $m - 1 \in \N$
के साथ $(m - 1) + 1$ है)। $\sigma$ के लिए: वह सम संख्याओं $0, 2, 4, \dots$
को $0, 1, 2, \dots$ पर और विषम संख्याओं $1, 3, 5, \dots$ को $-1, -2, -3,
\dots$ पर भेजता है। एकैकीयता: सम निवेश $\N$ में गिरते हैं ($\sigma(n) = n/2
\geq 0$) और विषम निवेश पूर्णतः ऋणात्मक पूर्णांकों में ($\sigma(n) = -(n+1)/2
\leq -1$), अतः टक्कर एक ही सम-विषम वर्ग के भीतर ही हो सकती है, जहाँ $\sigma$
पूर्णतः एकदिष्ट है ($n/2 = m/2$ या $(n+1)/2 = (m+1)/2$ से $n = m$ आ जाता
है)। आच्छादकता: $k \geq 0$ $\sigma(2k)$ है; $k \leq -1$ $\sigma(-2k - 1)$
है, जहाँ $-2k - 1 \geq 1$ विषम है। अतः $\N \approx \N^*$ और $\N \approx \Z$।

\textbf{5.} $C_0 = E \setminus g(F) \subseteq C$, अतः $x \notin C$ से $x
\notin C_0$ आता है, अर्थात् $x \in g(F)$: कोई $y \in F$ $g(y) = x$ को
संतुष्ट करता है। यदि $g(y') = x$ भी हो, तो $g$ की एकैकीयता से $y' = y$। अतः
$h$ की परिभाषा का दूसरा उपवाक्य एक अद्वितीय, सुपरिभाषित अवयव $g^{-1}(x)$ चुन
लेता है।

\textbf{6.} अग्र प्रतिबिंब सम्मिलन के साथ क्रम-विनिमेय हैं
(\cref{exo:b1:logic:8}~(1), पहले $f$ पर, फिर $g$ पर लगाकर):
\[
g\bigl(f(C)\bigr)
= g\Bigl(f\Bigl(\bigcup_{n \in \N} C_n\Bigr)\Bigr)
= \bigcup_{n \in \N} g\bigl(f(C_n)\bigr)
= \bigcup_{n \in \N} C_{n+1}
= \bigcup_{n \geq 1} C_n \subseteq C .
\]

\textbf{7.} मान लीजिए $E$ में $x \neq x'$। यदि दोनों $C$ में हों, तो $f$ की
एकैकीयता से $h(x) = f(x) \neq f(x') = h(x')$। यदि कोई भी $C$ में न हो, तो
$g(h(x)) = x \neq x' = g(h(x'))$, अतः $h(x) \neq h(x')$। यदि $x \in C$ और
$x' \notin C$ (मिश्रित स्थिति, नाम बदलने तक): मान लीजिए $h(x) = h(x')$,
अर्थात् $f(x) = g^{-1}(x')$। $g$ लगाने पर $x' = g(f(x)) \in g(f(C))$, और
प्रश्न 6 से $x' \in C$ मिलता है --- विरोधाभास। अतः सभी स्थितियों में $h(x)
\neq h(x')$: $h$ \omterm{def:b1:logic:inj}{एकैकी} है।

\textbf{8.} मान लीजिए $y \in F$। \emph{स्थिति 1: $g(y) \notin C$।} तब
$h(g(y)) = g^{-1}(g(y)) = y$: अवयव $g(y)$ एक \omterm{def:b1:logic:map}{पूर्वप्रतिबिंब} है। \emph{स्थिति
2: $g(y) \in C$}, मान लीजिए $g(y) \in C_n$। चूँकि $g(y) \in g(F)$, अतः $g(y)
\notin C_0 = E \setminus g(F)$, इसलिए $n \geq 1$ और $g(y) \in C_n =
g(f(C_{n-1}))$: ऐसा $x \in C_{n-1}$ है कि $g(y) = g(f(x))$। $g$ की एकैकीयता
से $y = f(x)$ मिलता है, और $x \in C_{n-1} \subseteq C$, अतः $h(x) = f(x) =
y$। दोनों स्थितियों में $y$ प्राप्त हो जाता है: $h$ \omterm{def:b1:logic:inj}{आच्छादक} है, अतः \omterm{def:b1:logic:inj}{एकैकी
आच्छादक}।

\textbf{9.} यदि $E \preceq F$ और $F \preceq E$, तो \omterm{def:b1:logic:inj}{एकैकी} \omterm{def:b1:logic:map}{प्रतिचित्रण} $f
\colon E \to F$ और $g \colon F \to E$ चुनिए; प्रश्न 5--8 एक एकैकी आच्छादन $h
\colon E \to F$ गढ़ देते हैं, अतः $E \approx F$। यह \omterm{def:b1:logic:statement}{कथन} इसलिए अतुच्छ है कि
दिए हुए दोनों \omterm{def:b1:logic:inj}{एकैकी} प्रतिचित्रणों का आपस में कोई संबंध नहीं है --- किसी का
\omterm{def:b1:logic:inj}{आच्छादक} होना आवश्यक नहीं, और $f$ तथा $g$ को भोलेपन से मिलाने वाला कोई सूत्र
\omterm{def:b1:logic:map}{प्रतिचित्रण} नहीं देता: सारी सामग्री $E$ के उस \omterm{thm:b1:logic:partition}{विभाजन} में है जो उसे क्षेत्र
$C$ (जहाँ $f$ की नकल की जाती है) और उसके पूरक (जहाँ $g$ को उलटी दिशा में
चलाया जाता है) में बाँटता है।

\textbf{10.} (क) अंतर्विष्टि $\intoo01 \to \intcc01$ \omterm{def:b1:logic:inj}{एकैकी} है; और $x \mapsto
\frac{x + 1}3$ $\intcc01$ को $\intcc{\frac13}{\frac23} \subseteq \intoo01$
में \omterm{def:b1:logic:inj}{एकैकी} रूप से भेजता है (वह अशून्य प्रवणता वाला ऐफ़ीन फलन है)। प्रश्न 9 से
$\intcc01 \approx \intoo01$ --- ऐसा एकैकी आच्छादन जिसे स्पष्ट रूप से लिखना
काफ़ी अप्रिय है। (ख) \emph{एकैकीयता।} मान लीजिए $2^p(2q + 1) = 2^{p'}(2q' +
1)$, जहाँ मान लीजिए $p \leq p'$। $2^p$ से भाग देने पर $2q + 1 = 2^{p' -
p}(2q' + 1)$। यदि $p' > p$, तो दायाँ पक्ष सम और बायाँ विषम होगा --- असंभव;
अतः $p = p'$, फिर $2q + 1 = 2q' + 1$ और $q = q'$। \emph{आच्छादकता।} प्रबल
आगमन से दिखाते हैं कि प्रत्येक पूर्णांक $m \geq 1$ $2^p(2q + 1)$ के रूप का
है। $m = 1$ के लिए: $p = q = 0$। मान लीजिए $m \geq 1$ और $\intint1m$ के सभी
पूर्णांकों के लिए दावा मान लीजिए। यदि $m + 1$ विषम है, तो $m + 1 = 2q + 1$,
जहाँ $p = 0$। यदि $m + 1$ सम है, तो $m + 1 = 2m'$, जहाँ $1 \leq m' \leq m$;
परिकल्पना से $m' = 2^p(2q + 1)$, अतः $m + 1 = 2^{p+1}(2q + 1)$। अतः
$\varphi(p, q) = 2^p(2q + 1) - 1$ प्रत्येक $n \in \N$ तक पहुँचता है, और
$\varphi$ एकैकी आच्छादन $\N \times \N \to \N$ है।

\textbf{11.} चूँकि $A$ अनंत है, $A \setminus \{\varphi(0), \dots, \varphi(n
- 1)\}$ कभी रिक्त नहीं होता, और $\N$ का न्यूनतम-अवयव गुणधर्म (जिससे
\cref{thm:b1:logic:induction} सिद्ध किया गया था) इस पुनरावर्ती परिभाषा को
वैध बना देता है। \emph{पूर्णतः वर्धमान:} $\varphi(n + 1)$ $A \setminus
\{\varphi(0), \dots, \varphi(n)\} \subseteq A \setminus \{\varphi(0), \dots,
\varphi(n - 1)\}$ का सदस्य है, जिसका न्यूनतम $\varphi(n)$ है; अतः $\varphi(n
+ 1) \geq \varphi(n)$, और समता संभव नहीं, जिससे $\varphi(n+1) > \varphi(n)$।
\emph{$\varphi(n) \geq n$:} आगमन से $\varphi(0) \geq 0$, और $\varphi(n + 1)
\geq \varphi(n) + 1 \geq n + 1$। \emph{एकैकीयता} कठोर एकदिष्टता से निकलती
है। \emph{$A$ पर आच्छादकता:} मान लीजिए कोई $a \in A$ कभी प्राप्त नहीं होता।
चूँकि $\varphi(a + 1) \geq a + 1 > a$, अतः $\varphi(n) > a$ वाले $n$ का
\omterm{def:b1:logic:sets}{समुच्चय} अरिक्त है; मान लीजिए $n$ उसका न्यूनतम अवयव है। प्रत्येक $k < n$ के
लिए $\varphi(k) \leq a$, अतः $\varphi(k) < a$ ($a$ प्राप्त नहीं होता)। तब
$a$ $A \setminus \{\varphi(0), \dots, \varphi(n - 1)\}$ में है और $a <
\varphi(n)$, जो $\varphi(n)$ को परिभाषित करने वाली लघुतमता का विरोध करता है।
अतः $\varphi$ एकैकी आच्छादन $\N \to A$ है, और $A$ गणनीय है।

\textbf{12.} मान लीजिए \omterm{def:b1:logic:inj}{एकैकी} \omterm{def:b1:logic:map}{प्रतिचित्रण} $f$ के द्वारा $E \preceq \N$; तब $E
\approx f(E)$ (प्रश्न 2)। यदि $f(E)$ परिमित है, तो $E$ परिमित है; यदि $f(E)$
अनंत है, तो प्रश्न 11 से $f(E) \approx \N$ मिलता है, अतः संक्रमकता (प्रश्न
1) से $E \approx \N$। विलोमतः, परिमित \omterm{def:b1:logic:sets}{समुच्चय} और \omterm{pb:b1:logic:1}{गणनीय समुच्चय} स्पष्टतः $\N$
में \omterm{def:b1:logic:inj}{एकैकी} रूप से समा जाते हैं। संक्षेप: $E \preceq \N$ और $\N \preceq E$ से
कैंटर--श्रोडर--बर्नस्टाइन द्वारा सीधे $E \approx \N$ मिल जाता है --- किसी
गणना-तर्क की आवश्यकता नहीं।

\textbf{13.} मान लीजिए $f \colon E \to \N$ और $g \colon F \to \N$ \omterm{def:b1:logic:inj}{एकैकी}
\omterm{def:b1:logic:map}{प्रतिचित्रण} हैं। तब $(x, y) \mapsto \varphi\bigl(f(x), g(y)\bigr)$ एक \omterm{def:b1:logic:inj}{एकैकी}
\omterm{def:b1:logic:map}{प्रतिचित्रण} $E \times F \to \N$ है: यदि प्रतिबिंब मेल खाएँ, तो $\varphi$
(प्रश्न 10) की एकैकीयता से $f(x) = f(x')$ और $g(y) = g(y')$ मिलते हैं, फिर
$x = x'$, $y = y'$। $\Z \times \N^*$ के लिए: दोनों गुणनखंड गणनीय हैं (प्रश्न
4), अतः $\Z \times \N^* \preceq \N$; वह अनंत है (उसमें $\{0\} \times \N^*$
है), अतः प्रश्न 12 से गणनीय है।

\textbf{14.} प्रत्येक परिमेय $r$ का अद्वितीय निरूपण $r = p/q$ है, जहाँ $p
\in \Z$, $q \in \N^*$ और भिन्न न्यूनतम पदों में है (अद्वितीयता
\cref{ch:b1:arith} में सिद्ध है; $r = 0$ के लिए $0/1$ लीजिए)। तब \omterm{def:b1:logic:map}{प्रतिचित्रण}
$r \mapsto (p, q)$ \omterm{def:b1:logic:inj}{एकैकी} है: युग्म $r = p/q$ को निर्धारित कर देता है। अतः
प्रश्न 13 से $\Q \preceq \Z \times \N^* \preceq \N$। चूँकि $\N \subseteq \Q$
से $\N \preceq \Q$ मिलता है, प्रश्न 12 (या सीधे कैंटर--श्रोडर--बर्नस्टाइन)
दिखाता है कि $\Q \approx \N$: परिमेय संख्याएँ गणनीय हैं।

\textbf{15.} प्रत्येक $n$ के लिए एक \omterm{def:b1:logic:inj}{एकैकी} \omterm{def:b1:logic:map}{प्रतिचित्रण} $f_n \colon E_n \to
\N$ नियत कीजिए। $x \in \bigcup_n E_n$ के लिए मान लीजिए $n(x)$ वह
\emph{न्यूनतम} $n$ है जिसके लिए $x \in E_n$, और $u(x) = \varphi\bigl(n(x),
f_{n(x)}(x)\bigr) \in \N$ रखिए। यदि $u(x) = u(x')$, तो $\varphi$ की एकैकीयता
से $n(x) = n(x') = n$ और $f_n(x) = f_n(x')$ मिलते हैं, अतः $f_n$ की एकैकीयता
से $x = x'$। अतः सम्मिलन $\N$ में \omterm{def:b1:logic:inj}{एकैकी} रूप से समा जाता है: वह अधिकतम गणनीय
है।

\textbf{16.} परिमित $F \subseteq \N$ ($\Psi(\emptyset) = 0$) के लिए $\Psi(F)
= \sum_{i \in F} 2^i$ रखिए। मान लीजिए $F \neq F'$ और $m$ वह बृहत्तम अवयव है
जिस पर वे भिन्न हैं, मान लीजिए $m \in F \setminus F'$ (आवश्यकता हो तो नाम
बदल लीजिए)। अवयव $> m$ या तो दोनों में हैं या किसी में नहीं, अतः वे दोनों
योगों में समान योगदान देते हैं; अवयवों $\leq m$ के योगदानों की तुलना कीजिए:
\[
\sum_{i \in F,\, i \leq m} 2^i \geq 2^m
> 2^m - 1 = \sum_{k=0}^{m-1} 2^k
\geq \sum_{i \in F',\, i \leq m} 2^i ,
\]
जहाँ \cref{exo:b1:logic:4} के गुणोत्तर योग का प्रयोग किया गया। अतः $\Psi(F)
\neq \Psi(F')$: $\Psi$ \omterm{def:b1:logic:inj}{एकैकी} है और $\N$ के परिमित उपसमुच्चयों का \omterm{def:b1:logic:sets}{समुच्चय}
अधिकतम गणनीय है; वह अनंत है (उसमें सभी एकल \omterm{def:b1:logic:sets}{समुच्चय} हैं), अतः गणनीय है।

\textbf{17.} $A \subseteq \N$ को उसके सूचक $\mathbf 1_A \colon \N \to
\{0,1\}$ पर भेजिए, जहाँ $n \in A$ होने पर $\mathbf 1_A(n) = 1$ और अन्यथा
$0$; और $u \in \{0,1\}^{\N}$ को $A_u = \{n \in \N : u(n) = 1\}$ पर भेजिए।
दोनों \omterm{def:b1:logic:map}{प्रतिचित्रण} परस्पर प्रतिलोम हैं: $A_{\mathbf 1_A} = A$ और $\mathbf
1_{A_u} = u$ (प्रत्येक $n$ पर मान जाँचिए)। \cref{thm:b1:logic:inverse} से
प्रत्येक एकैकी आच्छादन है: $\mathcal{P}(\N) \approx \{0,1\}^{\N}$।

\textbf{18.} प्रत्येक $n$ के लिए $d(n) = 1 - \Phi(n)(n) \neq \Phi(n)(n)$,
अतः अनुक्रम $d$ और $\Phi(n)$ सूचकांक $n$ पर भिन्न हैं: $d \neq \Phi(n)$। अतः
कोई $\Phi$ \omterm{def:b1:logic:inj}{आच्छादक} नहीं है, और प्रश्न 3 से कोई \omterm{def:b1:logic:inj}{एकैकी} \omterm{def:b1:logic:map}{प्रतिचित्रण}
$\{0,1\}^{\N} \to \N$ भी नहीं है: $\{0,1\}^{\N}$ अधिकतम गणनीय नहीं है।
प्रश्न 17 के शब्दकोश से, \omterm{def:b1:logic:map}{प्रतिचित्रण} $\Phi \colon \N \to \{0,1\}^{\N}$
वस्तुतः \omterm{def:b1:logic:map}{प्रतिचित्रण} $f \colon \N \to \mathcal{P}(\N)$ है, और $d$ \omterm{def:b1:logic:sets}{समुच्चय} $D
= \{n : n \notin f(n)\}$ के अनुरूप है (वस्तुतः $d(n) = 1 \iff \Phi(n)(n) = 0
\iff n \notin f(n)$): विकर्ण तर्क ही $E = \N$ के लिए \cref{exo:b1:logic:11}
की कैंटर वाली उपपत्ति \emph{है}।

\textbf{19.} $x_n = 0.d_1(n)\,d_2(n)\,d_3(n)\dots$ को उचित रूप में लिखिए और
$d_n(n) \neq 5$ होने पर $\delta_n = 5$, $d_n(n) = 5$ होने पर $\delta_n = 6$
परिभाषित कीजिए, फिर $x = 0.\delta_1\delta_2\delta_3\dots$। इस प्रसार में
केवल अंक $5$ और $6$ आते हैं, अतः वह $9$ की माला पर समाप्त नहीं होता: यह किसी
वास्तविक $x \in \intco01$ का उचित प्रसार है। प्रत्येक $n$ के लिए $x$ और
$x_n$ के $n$-वें अंक भिन्न हैं (रचना से $\delta_n \neq d_n(n)$); और चूँकि
उचित प्रसार अद्वितीय होते हैं, अतः $x \neq x_n$। इस प्रकार कोई अनुक्रम
$\intco01$ को नहीं भरता: फिर से प्रश्न 3 से, $\intco01$ अधिकतम गणनीय नहीं
है।

\textbf{20.} $\intco01 \subseteq \R$, अतः कोई \omterm{def:b1:logic:inj}{एकैकी} \omterm{def:b1:logic:map}{प्रतिचित्रण} $\R \to \N$
$\intco01$ पर सीमित होकर प्रश्न 19 का विरोध करेगा: $\R$ अगणनीय है। यदि $\R
\setminus \Q$ अधिकतम गणनीय होता, तो $\R = \Q \cup (\R \setminus \Q)$ दो
अधिकतम गणनीय समुच्चयों का सम्मिलन होता, अतः प्रश्न 15 से अधिकतम गणनीय होता
($E_0 = \Q$, और $n \geq 1$ के लिए $E_n = \R \setminus \Q$ लीजिए) ---
विरोधाभास। अतः अपरिमेय संख्याएँ अगणनीय हैं। ठीक-ठीक कहें तो: $\R$ के भीतर
परिमेय संख्याएँ एक \omterm{pb:b1:logic:1}{गणनीय समुच्चय} बनाती हैं जबकि उनका पूरक अगणनीय है; कोई
एकैकी आच्छादन $\R \setminus \Q$ को $\Q$ से कभी मिला नहीं सकता --- अपरिमेय
संख्याएँ परिमेय संख्याओं से कठोर अर्थ में ``अधिक'' हैं, यद्यपि दोनों अनंत
हैं और दोनों सघन हैं।

\textbf{21.} $p/q$ ($q \neq 0$ के साथ) $qX - p$ का मूल है, जो पूर्णांक
गुणांकों वाला अशून्य बहुपद है। $\sqrt 2$ $X^2 - 2$ का मूल है। $x = \sqrt 2 +
\sqrt 3$ के लिए: $x^2 = 5 + 2\sqrt 6$, अतः $x^2 - 5 = 2\sqrt 6$ और $(x^2 -
5)^2 = 24$, अर्थात्
\[
x^4 - 10x^2 + 1 = 0 :
\]
$\sqrt 2 + \sqrt 3$ $X^4 - 10X^2 + 1$ का मूल है।

\textbf{22.} $P = a_0 + a_1X + \dots + a_nX^n$ (घात $\leq n$, पूर्णांक
गुणांक) को $(a_0, \dots, a_n) \in \Z^{n+1}$ पर भेजिए: यह \omterm{def:b1:logic:inj}{एकैकी} है, क्योंकि
बहुपद अपने गुणांकों से निर्धारित होता है। $n$ पर आगमन से: $\Z^1 = \Z$ गणनीय
है (प्रश्न 4), और $\Z^{n+2} \approx \Z^{n+1} \times \Z$ प्रश्न 13 से अधिकतम
गणनीय है। अतः परिबद्ध घात वाले पूर्णांक बहुपदों का प्रत्येक \omterm{def:b1:logic:sets}{समुच्चय} अधिकतम
गणनीय है; वह अनंत है (उसमें अचर बहुपद हैं), अतः प्रश्न 12 से गणनीय है।

\textbf{23.} सभी पूर्णांक बहुपदों का \omterm{def:b1:logic:sets}{समुच्चय} $\bigcup_{n \in \N} \{P : \deg
P \leq n,\ P \text{ के पूर्णांक गुणांक हैं}\}$ है, जो गणनीय समुच्चयों का
गणनीय सम्मिलन है: प्रश्न 15 से अधिकतम गणनीय, और अनंत, अतः गणनीय।

\textbf{24.} प्रत्येक अशून्य पूर्णांक बहुपद $P$ के लिए मूलों का \omterm{def:b1:logic:sets}{समुच्चय} $R_P
= \{x \in \R : P(x) = 0\}$ परिमित है (अधिक से अधिक $\deg P$ अवयव, स्वीकृत)।
प्रश्न 23 से अशून्य पूर्णांक बहुपदों की गणना $P_0, P_1, P_2, \dots$ की जा
सकती है; तब $\mathcal{A} = \bigcup_{n \in \N} R_{P_n}$ परिमित (अतः अधिकतम
गणनीय) समुच्चयों का गणनीय सम्मिलन है: प्रश्न 15 से अधिकतम गणनीय। उसमें $\Q$
है (प्रश्न 21), अतः वह अनंत है: $\mathcal{A}$ गणनीय है।

\textbf{25.} यदि $\R \setminus \mathcal{A}$ अधिकतम गणनीय होता, तो $\R =
\mathcal{A} \cup (\R \setminus \mathcal{A})$ भी अधिकतम गणनीय होता (प्रश्न
15), जो प्रश्न 20 का विरोध करता है। अतः अबीजीय संख्याएँ विद्यमान हैं और वे
तो अगणनीय \omterm{def:b1:logic:sets}{समुच्चय} तक बनाती हैं, जबकि बीजीय संख्याएँ --- जिनमें पूर्णांकों से
करणी द्वारा बनी हर संख्या आती है --- $\R$ के भीतर मात्र एक गणनीय ढाँचा बनाती
हैं। पूरी रचना का सारांश: प्रश्न 1--3 तुलना की भाषा खड़ी करते हैं;
कैंटर--श्रोडर--बर्नस्टाइन (प्रश्न 5--9) एक चतुर एकैकी आच्छादन के बदले दो सरल
\omterm{def:b1:logic:inj}{एकैकी} प्रतिचित्रणों से समशक्तता सिद्ध करने देता है, और उसका प्रयोग $\intcc01
\approx \intoo01$ के लिए, $\Q$ के लिए तथा पूरे भाग V में हुआ; युग्मन एकैकी
आच्छादन (प्रश्न 10) ने गुणन और गणनीय सम्मिलनों (प्रश्न 13, 15) को शक्ति दी,
जिन्होंने बदले में $\Q$, पूर्णांक बहुपदों और $\mathcal{A}$ को शक्ति दी; और
विकर्ण तर्क (प्रश्न 18--19) ने वह एकमात्र कठोर असमिका $\N \prec \R$ दी जो
पूरी कहानी को अतुच्छ बनाती है। कैंटर का निष्कर्ष दार्शनिक दृष्टि से चकित
करने वाला है: उपपत्ति एक भी \omterm{pb:b1:logic:1}{अबीजीय संख्या} प्रस्तुत नहीं करती, फिर भी दिखा
देती है कि समशक्तता के अर्थ में \emph{लगभग हर} वास्तविक संख्या अबीजीय है।
किसी एक विशिष्ट \omterm{pb:b1:logic:1}{अबीजीय संख्या} का नाम लेने --- $\pi$ या $\eu$ --- के लिए
बिलकुल भिन्न गणित और कई दशकों का और परिश्रम चाहिए था।
\end{solution}
